% ----------------------------------------------------------------------------
%  Résumé / Abstract
% ----------------------------------------------------------------------------
\cleardoublepage
\phantomsection
\addcontentsline{toc}{chapter}{Résumé}
\markboth{Résumé}{Résumé}
\vspace*{1cm}
{\LARGE\bfseries\color{insatblue} Résumé\par}\vspace{2pt}\hrule\vspace{12pt}

\noindent
Ce projet compare trois manières de transmettre une image sur un canal sans fil bruité simulé, à budget
radio strictement identique. La première est la chaîne séparée de Shannon, telle qu'elle est déployée
aujourd'hui : codage de source, code convolutif, modulation QAM, décodage de Viterbi, le tout simulé bit
à bit. La deuxième est un codage conjoint source-canal profond (\emph{deep JSCC}) entraîné de bout en bout
sur une perte pixel. La troisième, proposée ici, reprend exactement la même architecture et ne modifie que
la fonction de perte : s'ajoutent à l'erreur quadratique un terme de caractéristiques et un terme de tâche
calculés à travers un classifieur gelé.

À architecture, bande et canal identiques, le système sémantique cède environ \SI{0.5}{\decibel} de PSNR
et gagne 10 à 14 points de précision chez un juge indépendant, en canal AWGN comme en canal de Rayleigh.
À \SI{0}{\decibel}, là où la chaîne séparée réelle tombe au niveau du hasard, les deux systèmes appris
restent exploitables : le codage conjoint supprime l'effet falaise. Deux extensions prolongent l'étude,
une transmission orientée tâche qui atteint une précision comparable avec 32 fois moins de bande, et un
renouvellement conditionnel des latents qui divise la bande par trois sur une scène vidéo statique. Les
limites sont énoncées explicitement : le canal est simulé, et sur de grandes images jamais vues à
l'entraînement la borne idéale de la chaîne séparée reste devant.

\vspace{10pt}
\noindent\textbf{Mots-clés :} communication sémantique, codage conjoint source-canal, apprentissage profond,
canal AWGN, évanouissement de Rayleigh, PSNR, SSIM, effet falaise, transmission orientée tâche.

\vspace{1.4cm}
{\LARGE\bfseries\color{insatblue} Abstract\par}\vspace{2pt}\hrule\vspace{12pt}

\noindent
This project compares three ways of sending an image over a simulated noisy wireless channel under a
strictly identical radio budget. The first is the separated Shannon chain as deployed today: source
coding, convolutional code, QAM modulation and Viterbi decoding, simulated bit by bit. The second is a
deep joint source-channel code (deep JSCC) trained end to end on a pixel loss. The third, proposed here,
keeps exactly the same architecture and changes only the loss function, adding a feature term and a task
term computed through a frozen classifier.

With identical architecture, bandwidth and channel, the semantic system gives up about \SI{0.5}{\decibel}
of PSNR and gains 10 to 14 accuracy points as measured by an independent judge, on AWGN as well as on
Rayleigh fading. At \SI{0}{\decibel}, where the real separated chain drops to chance level, both learned
systems remain usable: joint coding removes the cliff effect. Two extensions complete the study: a
task-oriented transmitter reaching comparable accuracy with 32 times less bandwidth, and conditional
replenishment of the latent tensor, which divides bandwidth by three on a static video scene. The
limitations are stated explicitly: the channel is simulated, and on large images never seen during
training the ideal bound of the separated chain remains ahead.

\vspace{10pt}
\noindent\textbf{Keywords:} semantic communication, joint source-channel coding, deep learning, AWGN
channel, Rayleigh fading, PSNR, SSIM, cliff effect, task-oriented transmission.
